\section{Síntesis y ampliación}
Esta clase gira en torno al inference-time compute: primero se deja que el prompt libere tokens, después se usa sampling/search para ampliar el solution space y, finalmente, se corrige el resultado mediante trace + self-correction. Las tres etapas se pueden combinar según se necesite, y en tareas de reasoning de alta calidad mejoran de manera notable la accuracy y la reliability.

También repasamos el compute scheduling (programación triangular), el mecanismo de coordinación entre self-consistency/search, el loop de control de stepwise verifier y Tree-of-Thought, así como la interpretación de los datos de benchmark, para formar un mapa de ruta completo que va del prompt al deployment.

\begin{importantbox}{Puntos clave para confirmar la implementación}
1. Dar seguimiento continuo, en la ingeniería, a las tres categorías de métricas: sampling, search y trace; 2. usar la curva de benchmark (por ejemplo, Arc AGI) como referencia para el budget gating; 3. incorporar el deployment SLO, el alert y el budget instrumentation al CI/CD pipeline, para garantizar que el inference-time compute no exceda un rango controlable.
\end{importantbox}

\begin{table}[H]
\centering
\caption{Patrones centrales del cómputo en tiempo de inferencia}
\begin{tabular}{p{0.28\textwidth}p{0.32\textwidth}p{0.32\textwidth}}
\toprule
Estrategia & Beneficio clave & Costo típico \\
\midrule
Ampliar el presupuesto de CoT & Permite que el modelo se acerque a la human-level accuracy en problemas difíciles & Más tokens, más latency, mayor sensibilidad en el diseño del prompt \\
Muestreo de múltiples candidatos y self-consistency & Permite recuperarse de mistakes; útil en escenarios sin verifier & Cada candidate implica su propio compute, y requiere un voter pool y una representación de consistencia \\
Tree-of-Thought y step verifier & Eliminar bad paths desde etapas tempranas, y mejorar la interpretabilidad en conjunto con el verifier & El número de branching se dispara, y hace falta diseñar una quality function \\
Trace feedback + self-correction & La puntuación line-by-line puede mejorar la accuracy de forma progresiva & Sin oracle se produce feedback drift; hace falta prompt guard rails \\
\bottomrule
\end{tabular}
\end{table}

\subsection{Resumen de la sección}
Las técnicas de tiempo de inferencia necesitan integrar el prompt, el search y el trace en un pipeline, y mediante la programación del presupuesto y el bucle de métricas hacer que cada módulo aporte su máximo beneficio en cada etapa.

\subsection{Lecturas adicionales}
\begin{itemize}
\item Wang et al., en «Self-Consistency improves chain of thought reasoning», explican el efecto de scaling del sampling + majority vote.
\item Zhou et al., en «Tree-of-Thought», proponen combinar el search con la step evaluation mediante una exploración en forma de árbol.
\item Artículos sobre trace feedback como «Trace-Enhanced LLM Safety» explican cómo el natural language trace ayuda al verifier.
\item La serie AlphaProof/AlphaGeometry/AlphaCode de DeepMind muestra los resultados de combinar inference-time RL + search + trace en tareas complejas de reasoning.
\end{itemize}

\end{document}
